\chapter{What is MobileCAD}

\section{In one sentence}

MobileCAD is \textbf{an Android front end for the FreeCAD engine}: it packs the desktop FreeCAD modeling
kernel, workbenches and Python ecosystem into a single APK, then rewrites the whole user interface in
Kotlin + Jetpack Compose — for touch screens, for tablets and for phones.

That means two things:

\begin{enumerate}
  \item What you get is not a “simplified” modeler. Parametric sketching, solid modeling, surfaces,
        assembly, drawings, FEM, CAM, BIM, meshes and point clouds all rest on FreeCAD's native C++
        implementation;
  \item The document format is fully interchangeable with FreeCAD. A part you modeled on the phone in a
        spare moment opens on the desktop and can be changed further, and the reverse is true as well.
\end{enumerate}

\begin{tip}[One engine, two interfaces]
Desktop FreeCAD draws its interface with Qt; MobileCAD draws it with Compose but does \emph{not} start
FreeCAD's Qt interface. It calls FreeCAD's App layer directly (\file{FreeCAD.so}, \file{Part.so},
\file{PartDesign.so}, …) and uses a “Qt emulation layer” written in Python so that FreeCAD's own task panels
(Pad parameters, BIM wall parameters, FEM constraints, …) still run — they are simply rendered as Compose
widgets. So every dialog you see has its logic in FreeCAD itself.
\end{tip}

\section{Capability map}

MobileCAD reproduces \textbf{980 of FreeCAD's 984 commands across 22 workbenches} (the parity matrix in the
source is \file{research/PARITY.md}). The only commands not ported are desktop window-management ones with no
phone counterpart. Grouped by purpose:

\begin{center}
\small
\begin{tabularx}{\linewidth}{@{}lX@{}}
\toprule
\textbf{Category} & \textbf{Workbenches and capabilities} \\
\midrule
Parametric modeling & \textbf{Sketch} (Sketcher): geometry, constraints, dimensions, construction lines,
B-splines; \textbf{Part Design}: Pad, Pocket, Revolution, Loft, Pipe, Helix, Hole, Fillet, Chamfer, Draft,
Thickness, Pattern, Mirrored, datum geometry \\
Solids and surfaces & \textbf{Part}: primitives, booleans, extrude/revolve/sweep/loft, offset, check,
measure; \textbf{Surface} \\
Assembly & \textbf{Assembly}: 12 joint types, grounding, solving, conflict diagnostics, BOM, snapshots,
simulation \\
Drawing & \textbf{Drawing} (TechDraw): pages and templates, projection groups, sections, detail views,
dimensions, symbols, export to PDF/SVG/DXF \\
Data and expressions & \textbf{Variables} (Spreadsheet): spreadsheets, aliases, cross-object expressions,
unit systems \\
Simulation and manufacturing & \textbf{FEM} (the CalculiX solver runs on the device); \textbf{CAM} (jobs,
tools, toolpaths, post-processing) \\
Architecture and exchange & \textbf{BIM} (walls, slabs, columns, windows, stairs, roofs, IFC import and
export); \textbf{Mesh}, \textbf{Points}, \textbf{OpenSCAD}, \textbf{Reverse Engineering}, \textbf{Robot},
\textbf{Inspection}, \textbf{Material} \\
Core tools & \textbf{Std}: file, edit, view, measure, macros, dependency graph, drawing export, preferences
and more \\
\bottomrule
\end{tabularx}
\end{center}

\begin{tip}[Is every one of those commands used often?]
No. Many of the 984 commands are synonyms or fine-grained operations of specialist workbenches. About
\textbf{40} are genuinely frequent in everyday modeling: the dozen or so drawing and constraint tools in
Sketch, Pad/Pocket/Fillet/Pattern in Part Design, and opening, saving and exporting files. This book covers
those in depth; for the rest, use Chapter 8 and the command reference in Appendix A.
\end{tip}

\section{How it relates to desktop FreeCAD}

\begin{center}
\small
\begin{tabularx}{\linewidth}{@{}lXX@{}}
\toprule
 & \textbf{Desktop FreeCAD} & \textbf{MobileCAD} \\
\midrule
Geometry kernel & Open CASCADE & The same (the same cross-compiled libraries) \\
Document format & \file{.FCStd} & The same; files open in either direction \\
Python console & Yes & No (but macros and Python features still run) \\
Task panels & Qt dialogs & The same Python panels, rendered as touch controls \\
Interface & Qt Widgets, two columns full of panels & Compose, one workspace plus floating panels \\
Input & Mouse and keyboard, shortcuts & Finger gestures + Numpad + on-screen keyboard; hardware keyboards
supported \\
Draft / CAM / FEM & Complete & Complete (the solvers ship inside the APK) \\
Desktop window management & Multiple windows, docking & None (not needed on mobile) \\
\bottomrule
\end{tabularx}
\end{center}

\begin{warn}[Keeping versions aligned]
Every MobileCAD version is bound to one FreeCAD commit. This manual corresponds to \textbf{FreeCAD main
e95a31c}; the subtitle in the app's top bar shows it, for example “Part · based on FreeCAD main e95a31c”. When
you open the same document in desktop FreeCAD, use a similar version, so that new features (some PartDesign
parameters, for instance) are not lost in an older one.
\end{warn}

\section{Requirements and performance expectations}

\begin{itemize}
  \item \textbf{System}: Android 8.0 or later on an \textbf{arm64} device (phones and tablets; Quest headsets
        running Android use a separate MobileCAD XR build).
  \item \textbf{Storage}: the APK is about 165 MB; once the first launch has unpacked the runtime, it occupies
        about 1 GB. Set aside 2 GB or more.
  \item \textbf{Memory}: 6 GB or more for smooth modeling; a tablet with 8--12 GB handles assemblies of
        several hundred features with ease.
  \item \textbf{Performance}: solid modeling is on a par with desktop FreeCAD. FEM solving is limited by phone
        thermals, so keep meshes to about a third of the size you would use on the desktop. Display of large
        assemblies is optimized for mobile (the viewport redraws on demand and uses off-screen caching).
\end{itemize}

\section{License, privacy and data}

MobileCAD is free software under the \textbf{GPL-3.0-or-later} and includes FreeCAD (LGPL), Open CASCADE
(LGPL), Qt (LGPL), Gmsh, CalculiX, IfcOpenShell and other components. All computation happens locally on the
device: \textbf{documents are never uploaded and no cloud is involved}. The only network access is the
“Check for updates” action you trigger yourself (it reads version information from GitHub Releases) and any
web page you open from BIM or Help. The full list of licenses is under
\menupath{Help → About MobileCAD → Open-source licenses}.
